\documentclass[11pt]{article}
\usepackage[margin=1in]{geometry}
\usepackage{parskip}
\begin{document}
\noindent Ook Lee\\
Department of Information System, Hanyang University\\
222 Wangsimni-ro, Seongdong-gu, Seoul 04763, Republic of Korea\\
ooklee@hanyang.ac.kr

\medskip
\noindent [Date]

\medskip
\noindent The Editor-in-Chief\\
IEEE Transactions on Reliability

\medskip
\noindent Dear Editor,

Please consider the enclosed manuscript, ``Condition-Based Maintenance of Large Language Models Under Knowledge Drift and Imperfect Repair,'' for publication in the IEEE Transactions on Reliability as a regular paper.

A deployed language model deteriorates. The world drifts away from its training data, and the usual remedy, a parameter update, can itself damage competence that users depend on. This is the structure of an imperfect-repair maintenance problem, and the manuscript treats it as one. Staleness is modeled as a shock-and-wear process driven by a marked point process, forgetting as an imperfect repair with a Fisher-information-weighted damage term, and the choice among repair mechanisms as an average-cost Markov decision process.

The main results are (i) a square-root law linking update cost, drift intensity, and the optimal maintenance interval; (ii) a comparative static showing that a repair leaving less residual damage should be applied more often; (iii) a nested control-limit result for a graduated repair menu; and (iv) a closed-form test for when an external retrieval store, which compensates for drift without repairing the parameters, is worth its upkeep. A measurement study on a public 70-million-parameter model, replicated over ten random seeds, supplies the model inputs. It also shows that the assumed ordering of repair actions has to be verified from measurements and cannot be read off a mechanism's name: an unreplayed low-rank adapter raised total error in every intervention.

We believe the paper suits the journal's readership because it extends imperfect-maintenance and condition-based maintenance theory to a new class of asset whose repair technology is gradient-based, and because it yields predictions that practitioners operating such assets can test on their own telemetry. The paper does not claim a closed-loop validation of the scheduling policy; Section~IX states this and the other limitations of the study.

The manuscript is original, has not been published, and is not under consideration elsewhere. The author declares no competing interests. Code, per-seed results, and pilot results for the measurement study are provided as supplementary material.

Sincerely,

\medskip
Ook Lee\\
Department of Information System, Hanyang University\\
ooklee@hanyang.ac.kr
\end{document}
